\section{讨论与局限}

四种方法的结果呈现出比较清楚的层次。短时 FFT 模板方法的 exact-match 为 0.265，macro-F1 为 0.244，接近四分类随机水平。这个结果说明，分帧、FFT、局部块和归一化相关可以构成完整的检测流程，但还不足以支撑稳定的四分类判别。中间案例也能解释这一点：\texttt{get}、\texttt{boy} 等样本的四个 FFT 相关分数常常非常接近，说明最大能量块和单边幅度谱模板不足以稳定区分 /g/、/b/、/d/ 这类短促爆破音。

基于 AI 的重实现方法把 exact-match 提升到 0.503，macro-F1 提升到 0.497。这个提升并不来自黑盒分类器，而是来自更合适的中间表示：log-mel 降维让谱形更平滑，以词首估计位置为中心的局部窗减少了整词元音的干扰，正负模板差分压低了类别共享能量，自相关补充了持续摩擦和短促爆破之间的时间结构。分标签结果也符合这个解释：/z/ 的 F1 达到 0.737，因为它的摩擦段更长、更稳定；/b/ 只有 0.359，说明较弱的爆破段和后接元音变化仍然会让它与 /g/、/d/ 混在一起。

MLP 和 CNN 用来观察学习型模型能否进一步超过基于 AI 的重实现方法。MLP 的 exact-match 为 0.485，低于基于 AI 的重实现方法，说明整词 MFCC 统计虽然稳定，但会把词尾元音和说话人差异一起平均进去。CNN 的 exact-match 达到 0.559，是当前最高；它在 /d/ 和 /z/ 上明显更强，说明扩大到 300 个词、6796 条语音后，卷积模型确实能从整张 log-mel 谱图中学到一部分局部模式。不过 epoch 消融也显示，dev macro-F1 最高的 30/40 epoch 设置并没有带来最高测试准确率，最终 25 个 epoch 的结果更好，因此 CNN 的优势仍然带有一定的训练方差和数据划分依赖。

因此，仅按最高分排序不足以概括这组结果。短时 FFT 模板方法分数最低，但最容易解释；基于 AI 的重实现方法不是最高分，却最能说明改进频谱表示和检测结构为什么有效；CNN 分数最高，但可解释性最弱。这个结果也提醒我们，在一热四分类设置下，label-wise accuracy 会被大量真阴性抬高，真正更值得关注的是 exact-match、macro-F1 和分标签 F1。

这项研究仍有几个限制。语音输入仍然是整词级音频，而不是人工切好的音素片段，因此短时 FFT 模板方法只能依靠滑窗搜索。/g/、/b/、/d/ 的爆破段很短，只要对齐误差稍大，模板证据就容易被后续元音淹没。当前的基于 AI 的重实现方法仍然属于模板法，没有做 phone-level alignment，也没有显式建模说话人和录音通道差异。当前 CNN 只是轻量对照，没有加入数据增强、预训练语音模型或强制对齐，因此不能代表深度语音模型的上限。
